\documentclass[11pt,a4paper]{article}
\usepackage[utf8]{inputenc}
\usepackage{amsmath}
\usepackage{amsfonts}
\usepackage{graphicx}
\usepackage{booktabs}
\usepackage{geometry}
\geometry{margin=1in}

% --- Project Metadata ---
\title{MATH3030: Mathematical Research Project\\ \Large Elliptic Curve Cryptography: An Algebraic Analysis of Secret Key Exchanges}
\author{Akari Oo \\ Student ID: 47510927}
\date{October 2026}

\begin{document}

\maketitle

\begin{abstract}
This research project explores the algebraic structures underpinning Elliptic Curve Cryptography (ECC) and demonstrates its application in secure key distribution frameworks. 
\end{abstract}

\newpage

\section{Introduction}
% Week 8/9 Focus: Writing for a generalist audience.
In our modern world, keeping digital data safe is incredibly important for everyday life, like using online banking or sending encrypted text messages on our phones. To protect this private information from hackers, computers use a system called public-key cryptography. This is a special method that allows two people to share a secret password over the public internet, even if they have never met before and are sitting on opposite sides of the world. 

For a long time, the most popular system for doing this was called RSA (Rivest--Shamir--Adleman). The math behind RSA is simple: it is very easy to multiply two massive prime numbers together, but it is nearly impossible for a computer to work backward and figure out what those original prime numbers were. However, because modern computers are getting faster every day, these numbers now have to be incredibly huge to stay safe from hackers. Managing these massive numbers slows down internet networks and drains the battery life of small mobile devices.

Elliptic Curve Cryptography (ECC) fixes this problem by switching to a completely different type of math. Instead of using massive prime numbers, it uses the geometric properties of curves. By applying the laws of mathematical groups that we study in abstract algebra, we can use these curves to create codes that are just as secure as traditional systems but use much smaller numbers. For example, a tiny 256-bit elliptic curve key provides the exact same security as a massive 3072-bit RSA key. This reduction in size makes data processing highly efficient, accelerates connection speeds, and saves battery power on mobile devices.

To explore how this math works, this report is broken down into three main sections. Section 2 will look at the mathematical foundations, showing how curves behave on a smooth graph and how they turn into a discrete grid of dots when we apply modular arithmetic. Section 3 will present my main project contribution: I will pick a simple curve equation, find its coordinates on a small number grid, and show a complete step-by-step example of two people sharing a secret key. Finally, Section 4 will discuss the security limits of these curves and look at how future quantum computers might impact them.


\section{Mathematical Exposition \& Foundations}
% Week 10 Focus: Rigorous abstract algebra.
\subsection{Elliptic Curves over Real Fields}
\subsection{Generalization over Finite Fields $\mathbb{F}_p$}
\subsection{The Group Addition Law \& Chord-and-Tangent Method}

\section{Algorithmic Simulation \& Key Exchange}
% Week 11 Focus: Custom curve math calculations.
\subsection{Parameter Selection \& Coordinate Mapping}
\subsection{Elliptic Curve Diffie-Hellman (ECDH) Simulation}

\section{Conclusion \& Future Work}
% Week 12 Focus: Wrap up and Post-Quantum cryptography.

\newpage
\appendix
\section{Project Management Reflection}
% Week 12/13 Focus: Version control logs and timeline evaluation.

\end{document}
% =========================================================================
% SECTION 2: MATHEMATICAL EXPOSITION & FOUNDATIONS
% =========================================================================

\subsection{Elliptic Curves over Real Fields}
An elliptic curve $E$ over the field of continuous real numbers $\mathbb{R}$ is defined algebraically by a cubic Weierstrass equation of the form:
\begin{equation}
    y^2 = x^3 + ax + b
    \label{eq:real_weierstrass}
\end{equation}
where the coefficients $a, b \in \mathbb{R}$ are fixed structural constants~\cite{hoffstein}. To ensure that the curve is smooth and well-behaved for cryptographic applications, we must prevent geometric anomalies such as cusps, self-intersections, or isolated nodes. This is achieved by enforcing a strict non-singularity condition, which dictates that the discriminant $\Delta$ of the cubic polynomial must be non-zero:
\begin{equation}
    \Delta = -16(4a^3 + 27b^2) \neq 0
    \label{eq:discriminant}
\end{equation}
Geometrically, a non-zero discriminant guarantees that the polynomial $x^3 + ax + b$ possesses three distinct roots (which may be real or complex), ensuring that a unique tangent line can be calculated at every single coordinate along the continuous curve path~\cite{washington}.

\subsection{Generalization over Finite Fields $\mathbb{F}_p$}
While real-valued continuous geometric curves are useful for visualization and modeling, practical digital cryptosystems cannot operate over infinite, continuous decimal spaces due to floating-point rounding errors. To achieve absolute mathematical precision, we restrict the domain of our coordinates to a finite field $\mathbb{F}_p$, representing the set of integers modulo a prime number $p$~\cite{washington}. Under this paradigm, the elliptic curve equation transforms into a modular congruence structure:
\begin{equation}
    y^2 \equiv x^3 + ax + b \pmod{p}
    \label{eq:modular_weierstrass}
\end{equation}
where the coefficients $a, b \in \mathbb{F}_p$ must satisfy the modified modular non-singularity check $4a^3 + 27b^2 \not\equiv 0 \pmod{p}$~\cite{hoffstein}. 

Shifting the domain to $\mathbb{F}_p$ alters the geometry fundamentally. Instead of a smooth, continuous line, the curve manifests as a finite, discrete grid of coordinates $(x, y) \in \mathbb{F}_p \times \mathbb{F}_p$ satisfying Equation~\ref{eq:modular_weierstrass}. The complete collection of these discrete solutions, denoted by $E(\mathbb{F}_p)$, is finite and always includes an idealized point at infinity, denoted as $\mathcal{O}$, which functions as the structural boundary element of the system.

\subsection{The Group Addition Law \& Chord-and-Tangent Method}
The finite set of discrete coordinates $E(\mathbb{F}_p)$ forms a commutative algebraic structure known as an abelian group under a specialized binary operation called point addition, denoted by the symbol $\oplus$~\cite{hoffstein}. To satisfy the formal axioms of group theory, the operation adheres to the following structural laws:
\begin{itemize}
    \item \textbf{Identity:} The point at infinity $\mathcal{O}$ serves as the additive identity element. For any point $P \in E(\mathbb{F}_p)$, $P \oplus \mathcal{O} = \mathcal{O} \oplus P = P$.
    \item \textbf{Inverse:} Every point $P = (x, y)$ possesses a unique additive inverse $-P = (x, -y \pmod{p})$, such that $P \oplus (-P) = \mathcal{O}$. Geometrically, the inverse is the reflection of the point across the central horizontal axis of symmetry.
    \item \textbf{Associativity \& Commutativity:} For any points $P, Q, R \in E(\mathbb{F}_p)$, the operations satisfy $(P \oplus Q) \oplus R = P \oplus (Q \oplus R)$ and $P \oplus Q = Q \oplus P$~\cite{washington}.
\end{itemize}

Geometrically over $\mathbb{R}$, adding two distinct points $P$ and $Q$ relies on the chord method: drawing a straight line through $P$ and $Q$ yields a third intersection point on the curve, $R'$, which is then reflected across the $x$-axis to discover the final sum $P \oplus Q = R$. If a point is added to itself ($P \oplus P$), a tangent line is constructed at $P$ to find the intersection point prior to reflection.

To evaluate these operations algebraically on our discrete modulo-$p$ space, let $P = (x_1, y_1)$ and $Q = (x_2, y_2)$ be elements of $E(\mathbb{F}_p)$. Assuming $P \neq -Q$, the coordinates of their group sum $R = P \oplus Q = (x_3, y_3)$ are computed using the following equations:
\begin{equation}
    x_3 \equiv \lambda^2 - x_1 - x_2 \pmod{p}
    \label{eq:add_x}
\end{equation}
\begin{equation}
    y_3 \equiv \lambda(x_1 - x_3) - y_1 \pmod{p}
    \label{eq:add_y}
\end{equation}

The modular slope parameter $\lambda$ is determined dynamically based on whether the operation is point addition or point doubling:
\begin{align}
    \lambda &\equiv (y_2 - y_1)(x_2 - x_1)^{-1} \pmod{p} \quad \text{when } P \neq Q \label{eq:slope_add} \\
    \lambda &\equiv (3x_1^2 + a)(2y_1)^{-1} \pmod{p} \quad \text{when } P = Q \label{eq:slope_double}
\end{align}
In Equations~\ref{eq:slope_add} and~\ref{eq:slope_double}, the exponent notation $(\cdot)^{-1}$ denotes the modular multiplicative inverse rather than a standard fractional denominator. If the modular inverse cannot be calculated because the denominator reduces to $0 \pmod{p}$, the geometric line is vertical, and the operation yields the identity point $\mathcal{O}$~\cite{hoffstein}.



% =========================================================================
% REFERENCES / BIBLIOGRAPHY
% =========================================================================
\newpage
\begin{thebibliography}{99}

\bibitem{hoffstein}
J.~Hoffstein, J.~Pipher, and J.~H.~Silverman, 
\textit{An Introduction to Mathematical Cryptography}, 2nd ed. 
New York, NY: Springer, 2014.

\bibitem{washington}
L.~C.~Washington, 
\textit{Elliptic Curves: Number Theory and Cryptography}, 2nd ed. 
Boca Raton, FL: Chapman \& Hall/CRC, 2008.

\end{thebibliography}
